% GENERATED FILE. Edit canonical lessons, then run npm run generate:books.
% canonical-source-hash: fnv1a64:f9f991339ed6b3cc
% canonical-lessons: GE-C285-betruegen, GE-C285-betrug, GE-C285-betrueger, GE-C285-faelschen, GE-C285-erpressen, GE-R285-review-278-285

\chapter{Der Betrug}
\label{ch:ge-der-betrug}
\hlchaptermodality{285}

\begin{chapteropening}
\textbf{By the end of this chapter:} \emph{I can report fraud: to cheat, fraud, a con man, to forge and to blackmail.}

Complete the last lesson of chapter 285: I can report fraud: to cheat, fraud, a con man, to forge and to blackmail.
\end{chapteropening}

\section[betrügen]{betrügen — to cheat, to deceive}

\label{lesson:GE-C285-betruegen}



\begin{quote}
Before the new one: say the German for the investigation, then the German for to arrest.
\end{quote}

\subsection*{You'll want to know: betrügen}
\textbf{betrügen} — \textquotedblleft{}to cheat, to deceive\textquotedblright{}.

\begin{quote}
Der Verkäufer betrog viele alte Leute.
\end{quote}

\begin{quote}
\emph{The salesman cheated many old people.}
\end{quote}

Strong verb: \emph{betrog}, \emph{hat betrogen}. Also to be unfaithful: \emph{Er betrügt seine Frau}.

\begin{cousinweb}
\emph{betrügen} is inherited from Old High German \emph{bitriogan}, which EtymDB traces to Proto-Germanic \emph{bidreuganą} \textquotedblleft{}to betray, to deceive\textquotedblright{}.
\end{cousinweb}

\begin{grammarlens}[title={What you've built}]
One more word for telling, explaining and describing at B1.
\end{grammarlens}

\subsection*{Your turn}
\begin{itemize}
\raggedright
  \item \emph{Say it:} \emph{betrügen}
  \item \emph{Say it:} \emph{betrügen}, once more
  \item \emph{Say it:} a sentence of your own with \emph{betrügen}
  \item \emph{Recall:} say the German for the investigation, then the German for to arrest, then say \emph{betrügen} again
\end{itemize}

\begin{tcolorbox}[breakable,colback=teal!4,colframe=teal!35!black,title={Before you move on}]
What does betrügen mean? (\textquotedblleft{}To cheat, to deceive\textquotedblright{}.) Say it once more.
\end{tcolorbox}
\section[der Betrug]{der Betrug — the fraud, the deception}

\label{lesson:GE-C285-betrug}



\begin{quote}
Before the new one: say the German for to arrest, then the German for to cheat.
\end{quote}

\subsection*{You'll want to know: der Betrug}
\textbf{der Betrug} — \textquotedblleft{}the fraud, the deception\textquotedblright{}.

\begin{quote}
Der Betrug kam nach einem Jahr ans Licht.
\end{quote}

\begin{quote}
\emph{The fraud came to light after a year.}
\end{quote}

No plural in most uses.

\begin{cousinweb}
\emph{Betrug} comes from Middle High German \emph{betroc}, Old High German \emph{bitrog} \textquotedblleft{}delusion\textquotedblright{} (EtymDB), from the same root as \emph{betrügen}.
\end{cousinweb}

\begin{grammarlens}[title={What you've built}]
One more word for telling, explaining and describing at B1.
\end{grammarlens}

\subsection*{Your turn}
\begin{itemize}
\raggedright
  \item \emph{Say it:} \emph{der Betrug}
  \item \emph{Say it:} \emph{der Betrug}, once more
  \item \emph{Say it:} a sentence of your own with \emph{der Betrug}
  \item \emph{Recall:} say the German for to arrest, then the German for to cheat, then say \emph{der Betrug} again
\end{itemize}

\begin{tcolorbox}[breakable,colback=teal!4,colframe=teal!35!black,title={Before you move on}]
What does der Betrug mean? (\textquotedblleft{}The fraud, the deception\textquotedblright{}.) Say it once more.
\end{tcolorbox}
\section[der Betrüger]{der Betrüger — the con man, the fraudster}

\label{lesson:GE-C285-betrueger}



\begin{quote}
Before the new one: say the German for to cheat, then the German for the fraud.
\end{quote}

\subsection*{You'll want to know: der Betrüger}
\textbf{der Betrüger} — \textquotedblleft{}the con man, the fraudster\textquotedblright{}.

\begin{quote}
Der Betrüger rief am Telefon an und wollte Geld.
\end{quote}

\begin{quote}
\emph{The con man phoned and wanted money.}
\end{quote}

Plural: \emph{die Betrüger}. Feminine \emph{die Betrügerin}.

\begin{cousinweb}
\emph{Betrüger} is built on \emph{betrügen} \textquotedblleft{}to cheat\textquotedblright{}, from Proto-Germanic \emph{bidreuganą} \textquotedblleft{}to deceive\textquotedblright{} (EtymDB), with the ending \emph{-er} for a person.
\end{cousinweb}

\begin{grammarlens}[title={What you've built}]
One more word for telling, explaining and describing at B1.
\end{grammarlens}

\subsection*{Your turn}
\begin{itemize}
\raggedright
  \item \emph{Say it:} \emph{der Betrüger}
  \item \emph{Say it:} \emph{der Betrüger}, once more
  \item \emph{Say it:} a sentence of your own with \emph{der Betrüger}
  \item \emph{Recall:} say the German for to cheat, then the German for the fraud, then say \emph{der Betrüger} again
\end{itemize}

\begin{tcolorbox}[breakable,colback=teal!4,colframe=teal!35!black,title={Before you move on}]
What does der Betrüger mean? (\textquotedblleft{}The con man, the fraudster\textquotedblright{}.) Say it once more.
\end{tcolorbox}
\section[fälschen]{fälschen — to forge, to fake}

\label{lesson:GE-C285-faelschen}



\begin{quote}
Before the new one: say the German for the fraud, then the German for the con man.
\end{quote}

\subsection*{You'll want to know: fälschen}
\textbf{fälschen} — \textquotedblleft{}to forge, to fake\textquotedblright{}.

\begin{quote}
Er fälschte einen Ausweis.
\end{quote}

\begin{quote}
\emph{He forged an identity card.}
\end{quote}

\emph{Gefälscht} — fake: \emph{ein gefälschter Pass}.

\begin{cousinweb}
\emph{fälschen} comes from Old High German \emph{felsken} (EtymDB), built on \emph{falsch} \textquotedblleft{}false\textquotedblright{}, which came from Latin \emph{falsus}, from \emph{fallere} \textquotedblleft{}to deceive\textquotedblright{} (Lewis \& Short).
\end{cousinweb}

\begin{grammarlens}[title={What you've built}]
One more word for telling, explaining and describing at B1.
\end{grammarlens}

\subsection*{Your turn}
\begin{itemize}
\raggedright
  \item \emph{Say it:} \emph{fälschen}
  \item \emph{Say it:} \emph{fälschen}, once more
  \item \emph{Say it:} a sentence of your own with \emph{fälschen}
  \item \emph{Recall:} say the German for the fraud, then the German for the con man, then say \emph{fälschen} again
\end{itemize}

\begin{tcolorbox}[breakable,colback=teal!4,colframe=teal!35!black,title={Before you move on}]
What does fälschen mean? (\textquotedblleft{}To forge, to fake\textquotedblright{}.) Say it once more.
\end{tcolorbox}
\section[erpressen]{erpressen — to blackmail, to extort}

\label{lesson:GE-C285-erpressen}



\begin{quote}
Before the new one: say the German for the con man, then the German for to forge.
\end{quote}

\subsection*{You'll want to know: erpressen}
\textbf{erpressen} — \textquotedblleft{}to blackmail, to extort\textquotedblright{}.

\begin{quote}
Ein Unbekannter erpresste die Firma.
\end{quote}

\begin{quote}
\emph{A stranger blackmailed the company.}
\end{quote}

Built on \emph{pressen} 'to press': to squeeze money out of someone. \emph{Die Erpressung} — blackmail.

\begin{cousinweb}
\emph{erpressen} is built on \emph{pressen} \textquotedblleft{}to press\textquotedblright{}, Old High German \emph{pressōn} (EtymDB), from Latin \emph{pressāre}, built on \emph{premere} \textquotedblleft{}to press\textquotedblright{} (Lewis \& Short).
\end{cousinweb}

\begin{grammarlens}[title={What you've built}]
One more word for telling, explaining and describing at B1.
\end{grammarlens}

\subsection*{Your turn}
\begin{itemize}
\raggedright
  \item \emph{Say it:} \emph{erpressen}
  \item \emph{Say it:} \emph{erpressen}, once more
  \item \emph{Say it:} a sentence of your own with \emph{erpressen}
  \item \emph{Recall:} say the German for the con man, then the German for to forge, then say \emph{erpressen} again
\end{itemize}

\begin{tcolorbox}[breakable,colback=teal!4,colframe=teal!35!black,title={Before you move on}]
What does erpressen mean? (\textquotedblleft{}To blackmail, to extort\textquotedblright{}.) Say it once more.
\end{tcolorbox}
\section[(dialogue)]{Words from chapters 278–285 — a review}

\label{lesson:GE-R285-review-278-285}



\begin{quote}
Say the German for to negotiate and the negotiation.
\end{quote}

\subsection*{Your turn}
Say these aloud:

\begin{itemize}
\raggedright
  \item to negotiate, the negotiation, the contract, to sign, the conference
  \item the event, to occur, to take place, to happen, the incident
  \item the crime, the criminal, the offender, the victim, the crime scene
  \item to break in, the burglary, the burglar, to pinch, the loot
  \item to hold up, the hold-up, to rob, to threaten, the weapon
  \item the witness, to testify, the statement, to observe, to recount
  \item the suspicion, suspicious, to investigate, the investigation, to arrest
  \item to cheat, the fraud, the con man, to forge, to blackmail
\end{itemize}

\begin{tcolorbox}[breakable,colback=teal!4,colframe=teal!35!black,title={Before you move on}]
To negotiate? (\textbf{verhandeln}.) The negotiation? (\textbf{die Verhandlung}.) The contract? (\textbf{der Vertrag}.) To sign? (\textbf{unterzeichnen}.) The conference? (\textbf{die Konferenz}.) The event? (\textbf{das Ereignis}.) To occur? (\textbf{ereignen}.) To take place? (\textbf{stattfinden}.) To happen? (\textbf{vorkommen}.) The incident? (\textbf{der Vorfall}.) The crime? (\textbf{das Verbrechen}.) The criminal? (\textbf{der Verbrecher}.) The offender? (\textbf{der Täter}.) The victim? (\textbf{das Opfer}.) The crime scene? (\textbf{der Tatort}.) To break in? (\textbf{einbrechen}.) The burglary? (\textbf{der Einbruch}.) The burglar? (\textbf{der Einbrecher}.) To pinch? (\textbf{klauen}.) The loot? (\textbf{die Beute}.) To hold up? (\textbf{überfallen}.) The hold-up? (\textbf{der Überfall}.) To rob? (\textbf{rauben}.) To threaten? (\textbf{bedrohen}.) The weapon? (\textbf{die Waffe}.) The witness? (\textbf{der Zeuge}.) To testify? (\textbf{aussagen}.) The statement? (\textbf{die Aussage}.) To observe? (\textbf{beobachten}.) To recount? (\textbf{schildern}.) The suspicion? (\textbf{der Verdacht}.) Suspicious? (\textbf{verdächtig}.) To investigate? (\textbf{ermitteln}.) The investigation? (\textbf{die Ermittlung}.) To arrest? (\textbf{festnehmen}.) To cheat? (\textbf{betrügen}.) The fraud? (\textbf{der Betrug}.) The con man? (\textbf{der Betrüger}.) To forge? (\textbf{fälschen}.) To blackmail? (\textbf{erpressen}.)
\end{tcolorbox}
